%\makeglossaries
\makenoidxglossaries

% ==================== DANH MỤC TỪ VIẾT TẮT ====================

\newglossaryentry{AI}{
    type=\acronymtype,
    name={AI},
    description={Trí tuệ nhân tạo (Artificial Intelligence)},
    first={Trí tuệ nhân tạo (Artificial Intelligence - AI)}
}

\newglossaryentry{API}{
    type=\acronymtype,
    name={API},
    description={Giao diện lập trình ứng dụng (Application Programming Interface)},
    first={Giao diện lập trình ứng dụng (Application Programming Interface - API)}
}

\newglossaryentry{BGM}{
    type=\acronymtype,
    name={BGM},
    description={Nhạc nền trò chơi (Background Music)},
    first={Nhạc nền trò chơi (Background Music - BGM)}
}

\newglossaryentry{FPS}{
    type=\acronymtype,
    name={FPS},
    description={Số khung hình trên giây (Frames Per Second)},
    first={Số khung hình trên giây (Frames Per Second - FPS)}
}

\newglossaryentry{FSM}{
    type=\acronymtype,
    name={FSM},
    description={Máy trạng thái hữu hạn (Finite State Machine)},
    first={Máy trạng thái hữu hạn (Finite State Machine - FSM)}
}

\newglossaryentry{HUD}{
    type=\acronymtype,
    name={HUD},
    description={Giao diện hiển thị thông tin trực quan trong màn chơi (Head-Up Display)},
    first={Giao diện hiển thị thông tin trực quan trong màn chơi (Head-Up Display - HUD)}
}

\newglossaryentry{IDE}{
    type=\acronymtype,
    name={IDE},
    description={Môi trường phát triển tích hợp (Integrated Development Environment)},
    first={Môi trường phát triển tích hợp (Integrated Development Environment - IDE)}
}

\newglossaryentry{JSON}{
    type=\acronymtype,
    name={JSON},
    description={Định dạng trao đổi dữ liệu gọn nhẹ (JavaScript Object Notation)},
    first={Định dạng trao đổi dữ liệu gọn nhẹ (JavaScript Object Notation - JSON)}
}

\newglossaryentry{LFS}{
    type=\acronymtype,
    name={LFS},
    description={Công cụ quản lý tệp tin dung lượng lớn của Git (Large File Storage)},
    first={Công cụ quản lý tệp tin dung lượng lớn của Git (Large File Storage - Git LFS)}
}

\newglossaryentry{NPC}{
    type=\acronymtype,
    name={NPC},
    description={Nhân vật không điều khiển được (Non-Player Character)},
    first={Nhân vật không điều khiển được (Non-Player Character - NPC)}
}

\newglossaryentry{PC}{
    type=\acronymtype,
    name={PC},
    description={Máy tính cá nhân (Personal Computer)},
    first={Máy tính cá nhân (Personal Computer - PC)}
}

\newglossaryentry{SFX}{
    type=\acronymtype,
    name={SFX},
    description={Hiệu ứng âm thanh (Sound Effects)},
    first={Hiệu ứng âm thanh (Sound Effects - SFX)}
}

\newglossaryentry{SO}{
    type=\acronymtype,
    name={SO},
    description={ScriptableObject - Cấu trúc lưu trữ dữ liệu tĩnh của động cơ Unity},
    first={ScriptableObject - SO}
}

\newglossaryentry{UI}{
    type=\acronymtype,
    name={UI},
    description={Giao diện người dùng (User Interface)},
    first={Giao diện người dùng (User Interface - UI)}
}

\newglossaryentry{UX}{
    type=\acronymtype,
    name={UX},
    description={Trải nghiệm người dùng (User Experience)},
    first={Trải nghiệm người dùng (User Experience - UX)}
}

\newglossaryentry{VFX}{
    type=\acronymtype,
    name={VFX},
    description={Hiệu ứng hình ảnh kỹ xảo (Visual Effects)},
    first={Hiệu ứng hình ảnh kỹ xảo (Visual Effects - VFX)}
}

\newglossaryentry{VRAM}{
    type=\acronymtype,
    name={VRAM},
    description={Bộ nhớ đồ họa ngẫu nhiên (Video Random Access Memory)},
    first={Bộ nhớ đồ họa ngẫu nhiên (Video Random Access Memory - VRAM)}
}


% ==================== DANH MỤC THUẬT NGỮ ====================

\newglossaryentry{Metroidvania}{
    name={Metroidvania},
    description={Thể loại trò chơi đi cảnh hành động kết hợp bản đồ liên thông phi tuyến tính và mở khóa kỹ năng để khám phá.}
}

\newglossaryentry{Sprite}{
    name={Sprite},
    description={Hình ảnh hai chiều (2D) tĩnh hoặc động được kết xuất và hiển thị trên màn hình trò chơi.}
}

\newglossaryentry{Collider}{
    name={Collider},
    description={Thành phần va chạm vật lý đại diện cho ranh giới hình học của thực thể trong động cơ vật lý.}
}

\newglossaryentry{Prefab}{
    name={Prefab},
    description={Hệ thống biểu mẫu định cấu hình sẵn của Unity, cho phép lưu trữ và tái sử dụng các GameObject cùng cấu trúc.}
}

\newglossaryentry{Animator}{
    name={Animator},
    description={Thành phần kiểm soát vòng đời và quy tắc chuyển đổi giữa các trạng thái hoạt ảnh của nhân vật trong Unity.}
}

\newglossaryentry{Culling}{
    name={Culling},
    description={Kỹ thuật tối ưu hóa bằng cách vô hiệu hóa hiển thị các đối tượng nằm ngoài tầm nhìn (Frustum/Camera View).}
}

\newglossaryentry{Lerp}{
    name={Lerp},
    description={Phương pháp nội suy tuyến tính (Linear Interpolation) dùng để mượt hóa quá trình chuyển đổi vị trí hoặc giá trị.}
}

\newglossaryentry{Vsync}{
    name={V-Sync},
    description={Cơ chế đồng bộ hóa dọc (Vertical Synchronization) giữa tốc độ khung hình trò chơi với tần số quét của màn hình.}
}

\newglossaryentry{MonoBehaviour}{
    name={MonoBehaviour},
    description={Lớp cơ sở mặc định trong Unity mà mọi tập lệnh lập trình C\# cần kế thừa để gắn kết vào các GameObject.}
}